% Options for packages loaded elsewhere
\PassOptionsToPackage{unicode}{hyperref}
\PassOptionsToPackage{hyphens}{url}
\documentclass[
]{ltjsarticle}
\usepackage{xcolor}
\usepackage{amsmath,amssymb}
\setcounter{secnumdepth}{-\maxdimen} % remove section numbering
\usepackage{iftex}
\ifPDFTeX
  \usepackage[T1]{fontenc}
  \usepackage[utf8]{inputenc}
  \usepackage{textcomp} % provide euro and other symbols
\else % if luatex or xetex
  \usepackage{unicode-math} % this also loads fontspec
  \defaultfontfeatures{Scale=MatchLowercase}
  \defaultfontfeatures[\rmfamily]{Ligatures=TeX,Scale=1}
\fi
\usepackage{lmodern}
\ifPDFTeX\else
  % xetex/luatex font selection
\fi
% Use upquote if available, for straight quotes in verbatim environments
\IfFileExists{upquote.sty}{\usepackage{upquote}}{}
\IfFileExists{microtype.sty}{% use microtype if available
  \usepackage[]{microtype}
  \UseMicrotypeSet[protrusion]{basicmath} % disable protrusion for tt fonts
}{}
\makeatletter
\@ifundefined{KOMAClassName}{% if non-KOMA class
  \IfFileExists{parskip.sty}{%
    \usepackage{parskip}
  }{% else
    \setlength{\parindent}{0pt}
    \setlength{\parskip}{6pt plus 2pt minus 1pt}}
}{% if KOMA class
  \KOMAoptions{parskip=half}}
\makeatother
\usepackage{longtable,booktabs,array}
\newcounter{none} % for unnumbered tables
\usepackage{calc} % for calculating minipage widths
% Correct order of tables after \paragraph or \subparagraph
\usepackage{etoolbox}
\makeatletter
\patchcmd\longtable{\par}{\if@noskipsec\mbox{}\fi\par}{}{}
\makeatother
% Allow footnotes in longtable head/foot
\IfFileExists{footnotehyper.sty}{\usepackage{footnotehyper}}{\usepackage{footnote}}
\makesavenoteenv{longtable}
\setlength{\emergencystretch}{3em} % prevent overfull lines
\providecommand{\tightlist}{%
  \setlength{\itemsep}{0pt}\setlength{\parskip}{0pt}}
\usepackage{bookmark}
\IfFileExists{xurl.sty}{\usepackage{xurl}}{} % add URL line breaks if available
\urlstyle{same}
\hypersetup{
  hidelinks,
  pdfcreator={LaTeX via pandoc}}

\author{}
\date{}

\begin{document}

\section{あるモデル粒子をめぐる思考実験------位相面積の床から次元の許容条件へ}\label{ux3042ux308bux30e2ux30c7ux30ebux7c92ux5b50ux3092ux3081ux3050ux308bux601dux8003ux5b9fux9a13ux4f4dux76f8ux9762ux7a4dux306eux5e8aux304bux3089ux6b21ux5143ux306eux8a31ux5bb9ux6761ux4ef6ux3078}

\textbf{著者}：木原 範昭（Noriaki Kihara） \textbf{所属}：WF System Co.,
Ltd.~/ 大阪大学基礎工学部（卒業）
\textbf{ORCID}：\href{https://orcid.org/0009-0004-6753-4020}{0009-0004-6753-4020}
\textbf{版}：v1（ステルス・ドラフト） \textbf{日付}：2026 年 6 月
\textbf{ライセンス}：CC BY 4.0 \textbf{Concept
DOI}：\href{https://doi.org/10.5281/zenodo.20528511}{10.5281/zenodo.20528511}
\textbf{Version DOI
(v1.0)}：\href{https://doi.org/10.5281/zenodo.20528512}{10.5281/zenodo.20528512}

\begin{center}\rule{0.5\linewidth}{0.5pt}\end{center}

\subsection{本稿の性格}\label{ux672cux7a3fux306eux6027ux683c}

\textbf{本稿は観察論文であり、思考実験の形式を採る。証明論文・主張論文ではない。}

本稿は標準量子論を修正しない。新しい物理理論を提案しない。新しい数学的定理を証明しない。観測可能な予言を一切変更しない。本稿が行うのは、\textbf{一つのモデル粒子を仮定する思考実験}を通じて、すでに標準的に確立された二つの数学------Robertson
の不確定性の床（位相空間における面積ゼロ状態の不在）と、ノルム付き多元体が次元
1, 2, 4, 8 でのみ存在するという Hurwitz--Frobenius
の系列------を、\textbf{位相空間の面積という一本の視点で並べ直して観察する}ことに尽きる。

本稿で扱う「粒子」は、特定の実在粒子（光子・電子等）ではなく、\textbf{位相面を一枚内在させた抽象的なモデル粒子}である。具体名を出さないのは、思考実験の一般性を保つためであり、特定の統計や質量・スピンに依存した主張を行わないことを明示するためである。

本稿は以下を主張しない：

\begin{itemize}
\tightlist
\item
  標準量子論・場の量子論・特殊相対論の数学的予言の変更
\item
  時空次元が 3+1
  であることの物理的導出（本稿は次元の数学的許容条件のみを観察する）
\item
  計量符号・ミンコフスキー構造・虚数単位の導出
\item
  新しい物理定数・断面積・崩壊率・新粒子の予言
\item
  新しい数学的定理の証明
\item
  宇宙の起源に関する形而上学的主張
\end{itemize}

評価・解釈は読者に委ねる。

\begin{center}\rule{0.5\linewidth}{0.5pt}\end{center}

\subsection{要旨}\label{ux8981ux65e8}

一つのモデル粒子を仮定し、それが位相平面 \((q,p)\)
の上にどのように存在しうるかを思考実験として追う。出発点は二つの既知の事実である。第一に、Robertson
の不確定性関係は、位相空間における面積ゼロの点状状態（\(\Delta q=\Delta p=0\)）を禁じる。最小のシンプレクティック面積
\(\sqrt{\det\Sigma}=\tfrac12\)（\(\hbar=1\)
単位）が床として残る。第二に、ノルムが積で保たれる多元体（平方和の積公式が成り立つ代数）は、実数・複素数・四元数・八元数の四つに限られ、その次元は
1, 2, 4, 8 である（Hurwitz）。

この二つを位相面積の視点で並べる。ただし以下の (ii)(iii) と §6
の読み替えは標準理論からの導出ではなく、振動自由度を時空四軸へ接続して読む\textbf{著者の対応仮説}である（本文「本稿後半について」で三層に分けて明示する）。(i)
モデル粒子は点に潰れられず、必ず最小面積 \(\tfrac12\)
の広がりを持つ。等方最小不確定状態の各軸標準偏差は \(\tfrac{1}{\sqrt2}\)
であり、本稿の表記 \(\pm\tfrac12\) は各軸の幅ではなく床
\(\tfrac12\)（ゼロ点の半端）そのものを指す。(ii)
位相面の二軸を、ともに振動数の次元へ無次元化して \((\nu_1,\nu_2)\)
と取ると、形式的には換算定数を陽に挟まず積 \(\nu_1\cdot\nu_2\)
と書ける（自然単位と写像の選択であり、\(\hbar\)
の物理的役割が消えるわけではない。なお振幅・振動数は標準的な正準共役対ではない）。比
\(k=\nu_2/\nu_1\) がスクイーズ・パラメータとして振る舞い、\(k=1\)
が等方、\(k\neq1\) が非等方に対応する。(iii)
運動量とエネルギーを空間方向・時間方向の振動数の従属量として読むと、四本の振動数の従属量を並べた\textbf{形式的な四成分量}が得られる。これが物理的な四元運動量
\((E,\boldsymbol{p})\) になるには、ミンコフスキー計量の別途導入が要る。

最後に、次元の許容条件について観察する。位相面積の床は連続的な制約であり、それ自体では次元を整数に固定しない。次元を四という鋭い整数値に定めるのは、多元体の存在という、床よりも厳しい離散的条件である。篩にかける条件は本稿が後付けで選んだものではなく、標準量子論の構造に現れる対称性------ノルム保存（ユニタリ性）、除法可能（可逆発展）、回転の非可換（角運動量代数。次元
1, 2 を排除）、合成の結合（演算子合成。次元 8
を排除）------を、振動数四軸の合成代数の上へ移して読んだものである。この読みのもとで、これらを同時に満たす最小の次元が四（四元数
\(\mathbb{H}\)）となる（§6.3 に次元 1〜6
の一覧表として示す）。なお本稿は、状態空間が複素数体（\(\mathbb{C}\)）であることは標準どおり継承し、四元数と結びつくのは状態のスカラー体ではなく時空四軸の合成代数である。連続の床（点になれない）と離散の環（四でのみ閉じる）という二つの独立な制約が同じ次元四に合流する、という符合を記録する。なお、系がなぜ四まで達し四で止まるか（次元の安定化）については、本稿は証明せず、§7
で「次元数もまた不確定性を負う」という仮定のもとでの推測として観測するに留める（類似する先行研究
{[}12{]}{[}13{]} は文脈であって論拠ではない）。

本稿は新しい主張を行わず、これらの既知の構造を一つの思考実験の中で並置観察するに留まる。符号（ミンコフスキー計量）の導入は本稿では扱わず、今後の課題とする。

\begin{center}\rule{0.5\linewidth}{0.5pt}\end{center}

\subsection{§1
思考実験の設定：一つのモデル粒子}\label{ux601dux8003ux5b9fux9a13ux306eux8a2dux5b9aux4e00ux3064ux306eux30e2ux30c7ux30ebux7c92ux5b50}

\subsubsection{1.1
何を仮定し、何を仮定しないか}\label{ux4f55ux3092ux4eeeux5b9aux3057ux4f55ux3092ux4eeeux5b9aux3057ux306aux3044ux304b}

一つのモデル粒子を考える。この粒子について、本稿が仮定するのは次の一点のみである------それは振動する自由度を持つ。質量・電荷・スピン・統計性、および背景時空（外から与えられた座標軸や時計）は仮定しない。具体的な実在粒子に同定することも避ける。光子や電子といった名を出すと、「なぜその粒子か」「ボソンだからか、質量ゼロだからか」という、本稿が扱わない問いを呼び込むためである。本稿の粒子はあくまで、位相面を一枚持つ抽象的な振動の単位である。

この最小の仮定から出発して、「振動する単位が存在するとは、位相空間においてどういうことか」を順に追う。各段で用いる数学はすべて標準的に確立されたものであり、本稿はそれらを引用し、一つの視点で並べるだけである。

\subsubsection{\texorpdfstring{1.2 舞台：位相平面
\((q,p)\)}{1.2 舞台：位相平面 (q,p)}}\label{ux821eux53f0ux4f4dux76f8ux5e73ux9762-qp}

観測者がこの粒子について量を読む舞台を、共役対 \((q,p)\)
の張る位相平面に取る。位相平面で意味を持つ不変量は、長さでも体積でもなく\textbf{面積}（作用次元）である。\(q\)
と \(p\)
が一組の共役ペアであり、面積を保つ変換がシンプレクティック変換だからである。この設定は同著者の観察論文
{[}10{]}
と共通であり、本稿はその枠組みの上で思考実験を進める。以下、面積と言うときは一貫してシンプレクティック面積
\(\sqrt{\det\Sigma}\) を指す（\(\Sigma\) は共分散行列）。

\begin{center}\rule{0.5\linewidth}{0.5pt}\end{center}

\subsection{\texorpdfstring{§2 点になれない------床 \(\tfrac12\)
と位相面の内在}{§2 点になれない------床 \textbackslash tfrac12 と位相面の内在}}\label{ux70b9ux306bux306aux308cux306aux3044ux5e8a-tfrac12-ux3068ux4f4dux76f8ux9762ux306eux5185ux5728}

\subsubsection{2.1 Robertson
の床が禁じるもの}\label{robertson-ux306eux5e8aux304cux7981ux3058ux308bux3082ux306e}

Robertson の不確定性関係（Robertson 1929 {[}2{]}）は

\[\Delta A\cdot\Delta B \ge \tfrac12\bigl|\langle[A,B]\rangle\bigr|\]

であり、\(A=q,\ B=p,\ [q,p]=i\)（\(\hbar=1\)）を入れると
\(\Delta q\cdot\Delta p\ge\tfrac12\)
を与える。より正確には、\(\mathrm{Sp}(2,\mathbb{R})\)
不変なシンプレクティック面積

\[\sqrt{\det\Sigma}=\sqrt{(\Delta q)^2(\Delta p)^2-\mathrm{Cov}(q,p)^2}\ \ge\ \tfrac12\]

が床を持つ（de Gosson の量子ブロブ {[}4{]}）。

ここで正確に述べる。禁じられるのは、位相空間における\textbf{面積ゼロの状態}（\(\Delta q=\Delta p=0\)
を満たす点状状態）であって、平均値としての原点
\((\langle q\rangle,\langle p\rangle)=(0,0)\)
ではない。後者は通常の状態として取れる。したがって本稿の観察は「面積ゼロの点状状態は存在しない」であり、「座標原点が存在しない」ではない。

\subsubsection{2.2
振動が存在するには面が要る}\label{ux632fux52d5ux304cux5b58ux5728ux3059ux308bux306bux306fux9762ux304cux8981ux308b}

モデル粒子が振動する、とは位相面の上で位相が進むことである。位相の進行（回転）には、回転すべき\textbf{面}が必要であり、面は最低二本の軸（共役対）が張る。一本の軸の上では回転＝振動を定義できない。

ここで思考実験の最初の分岐点に注意する。もしモデル粒子の振動を支える面を粒子の\textbf{外部}（背景時空）に探しにいくと、背景時空という未仮定の構造を密輸入することになる。本稿はこれを避け、位相面を粒子に\textbf{内在}するものと見る。場の各モードが調和振動子の構造（共役対
\(q,p\)）を内蔵していることは標準的な事実であり、本稿はこの内在する面の上で振動を読む。すなわち、振動を回す相手は外に要らず、面はモデル粒子の内側に初めから在ると仮定する。

\subsubsection{\texorpdfstring{2.3 床 \(\tfrac12\) とそのゼロ点表記
\(\pm\tfrac12\)}{2.3 床 \textbackslash tfrac12 とそのゼロ点表記 \textbackslash pm\textbackslash tfrac12}}\label{ux5e8a-tfrac12-ux3068ux305dux306eux30bcux30edux70b9ux8868ux8a18-pmtfrac12}

相関ゼロ（\(\mathrm{Cov}(q,p)=0\)）の等方な最小不確定状態では、各軸の標準偏差は
\(\Delta q=\Delta p=\tfrac{1}{\sqrt2}\) であり、その積が床
\(\Delta q\cdot\Delta p=\tfrac12\) になる。

ここで用語を正確に定義する。本稿が以降で \(\pm\tfrac12\)
と書くのは、\textbf{各軸の標準偏差ではない}（それは
\(\tfrac{1}{\sqrt2}\approx0.707\) であって \(\tfrac12\)
ではない）。\(\pm\tfrac12\) が指すのは、\textbf{面積の床 \(\tfrac12\)
そのもの}------半古典量子化 \(\oint p\,dq=(n+\tfrac12)h\) のゼロ点
\(\tfrac12\) にあたる、絞りきれない半端な量------である。\(\pm\)
は、この半端な床が共役二方向の双方に効くことを表す符号であって、軸ごとの幅の数値ではない。したがって
\(0.5\times0.5\)
のような積を取る対象ではなく、床という一つの量を二方向に帰属させた表記である。

重要なのは、この床 \(\tfrac12\)
が外部の時計や物差しで測った「振動数の不確定さ」ではない点である。それは、モデル粒子が点に潰れられないという内在的な広がり------位相面の最小の面積------そのものである。

\begin{center}\rule{0.5\linewidth}{0.5pt}\end{center}

\subsection{本稿後半について------既知の数学と、著者の対応仮説の境界}\label{ux672cux7a3fux5f8cux534aux306bux3064ux3044ux3066ux65e2ux77e5ux306eux6570ux5b66ux3068ux8457ux8005ux306eux5bfeux5fdcux4eeeux8aacux306eux5883ux754c}

ここから本稿の性格が分岐するため、議論を三つの層に分けて、それぞれの地位を最初に明示する。以降（§3〜§7）を読む際の前提である。

\textbf{第一層：確立した数学（観察の対象）。} §2 で用いた Robertson
の不確定性関係と、§6 で参照する Hurwitz--Frobenius
型のノルム付き多元体の次元制限（1, 2, 4,
8）は、いずれも既知の確立された数学的事実である。本稿はこれらを証明せず、引用して並置する。

\textbf{第二層：著者の対応仮説（本稿独自の読み替え）。}
これらの数学を「モデル粒子に内在する振動自由度」から「時空四軸」へ接続して読む部分------位相面の二軸を振幅と振動数に取ること（§3）、両軸を同種の無次元量に揃えること（§4）、四本の振動数に四元運動量を対応させること（§5）、四軸の合成代数に量子論の対称性を移して四元数を選ぶこと（§6.3）------は、\textbf{標準理論からの導出ではない}。これは著者の構造的直観にもとづく\textbf{作業仮説（対応仮定）}であり、その妥当性を本稿は証明しない。とくに、振幅と振動数が標準的な正準共役対であること、時空四軸の合成代数がノルム付き除法多元体でなければならないことは、いずれも本稿が課す仮定であって、標準量子論が要求するものではない。

\textbf{第三層：さらに推測的な部分。} §7
の「次元数もまた不確定性を負う」「四で安定する」は、第二層よりさらに踏み込んだ、著者の非導出的な推測である。

以上を踏まえ、以降の「四本の振動数」「四元数」「時空四軸」への接続は、証明ではなく\textbf{対応仮説}として読まれたい。本稿が記録するのは、出自の異なる二つの数学（位相面積の床と多元体の離散次元系列）が、モデル粒子の振動自由度という一つの視点の下では同じ次元四を指すように見える、という構造的符合である。本稿はこの符合を導出と称さず、また現実の時空が四次元であることを導くものでもない。この立場を省略すると本稿の主旨（何を観察し、何を仮説として置いているか）が見えなくなるため、削らずにここで明示する。

\begin{center}\rule{0.5\linewidth}{0.5pt}\end{center}

\subsection{§3
位相面の二軸------振幅と振動数、従属ラベルとしての時間}\label{ux4f4dux76f8ux9762ux306eux4e8cux8ef8ux632fux5e45ux3068ux632fux52d5ux6570ux5f93ux5c5eux30e9ux30d9ux30ebux3068ux3057ux3066ux306eux6642ux9593}

\subsubsection{3.1
共役軸のラベル選択}\label{ux5171ux5f79ux8ef8ux306eux30e9ux30d9ux30ebux9078ux629e}

位相面の二軸に、どのラベルを貼るか。慣習的には \((q,p)=\)（位置,
運動量）あるいは（時間,
振動数）と取る。しかし本稿の思考実験では、背景時空（外部の時計・物差し）を仮定していない。外部時間を共役軸に据えると、未仮定の構造を密輸入することになる。

そこで、モデル粒子に内在する量だけで面が閉じるラベルを選ぶ。一方の軸を\textbf{振幅}（振れの大きさ）、他方を\textbf{振動数}（位相が回る速さ）と取る。振幅は粒子に内在し、振動数も内在する。

ここを正確に断る。振幅と振動数は、標準的な意味での正準共役対------\([q,p]=i\)
を満たす対------ではない。調和振動子では振動数はハミルトニアンのパラメータであり、振幅は状態の大きさであって、両者が正準交換関係を持つわけではない（時間--帯域の
Gabor 不確定性はあるが、これは Robertson
型の正準交換関係とは性格が異なる）。本稿はこの二軸を、振動モードの状態を記述する\textbf{有効二変数}として形式的・観察的に導入するのであって、\((振幅, 振動数)\)
が正準対であると主張するものではない。Robertson
の床をこの面に移して読むこと自体が、「本稿後半について」で述べた対応仮説に属する。時間は、この面の上で位相がどれだけ進んだかを外から読むための\textbf{従属ラベル}として扱える。「単位時間あたり何回」という言い方は、面の回転を外部基準へ射影して初めて成り立つ。したがって本稿では、\textbf{面の主軸を振幅と振動数に取り、時間は従属量として扱う}------この取り方を仮定として置く。これは「時間が主軸ではありえない」という排他的な主張ではなく、振動数を主軸とする取り方が本稿の枠組みと矛盾しない、という意味である（時間を主軸に据える別の取り方を排除したわけではない）。

\subsubsection{3.2
振幅と振動数の反相関}\label{ux632fux5e45ux3068ux632fux52d5ux6570ux306eux53cdux76f8ux95a2}

振幅軸と振動数軸が面積 \(\tfrac12\)
を張るとき、片方を絞れば他方が伸びる。振動数を大きく（速く細かく振れる側に）寄せれば振幅側の広がりが詰まり、振動数を小さく（極限まで波長が長い側に）寄せれば振幅側が伸びる。面積は不変のまま、形だけが伸縮する。これは標準的な面積保存変換（スクイーズ）であり、同種の変換は同著者
{[}10{]} でも扱うが、本稿の以下の議論はそれに依存しない。

\begin{center}\rule{0.5\linewidth}{0.5pt}\end{center}

\subsection{\texorpdfstring{§4 無次元化と比
\(k\)------等方と分化}{§4 無次元化と比 k------等方と分化}}\label{ux7121ux6b21ux5143ux5316ux3068ux6bd4-kux7b49ux65b9ux3068ux5206ux5316}

\subsubsection{4.1
両軸を振動数に揃える}\label{ux4e21ux8ef8ux3092ux632fux52d5ux6570ux306bux63c3ux3048ux308b}

振幅軸と振動数軸を別種の次元のまま扱うと、面積を作るのに換算定数（\(\hbar\)）を挟むことになる。そこで両軸をともに振動数の次元に無次元化し、\((\nu_1,\nu_2)\)
と書く。この無次元化は、モデル粒子の内在する位相平面だけで面積を表現するためのものであり、背景時空由来の次元換算を介在させないための便宜である。面積は純粋な積

\[\text{面積}\ \propto\ \nu_1\cdot\nu_2\]

となり、形式的には換算定数を陽に挟まずに書ける。床は無次元の
\(\tfrac12\) である。ここも正確に述べる。これは自然単位系 \(\hbar=1\)
を採用したうえで、二軸を同種の無次元変数へ写す\textbf{写像の選択}であって、標準物理において
\(\hbar\)
が物理的役割を失うことを意味しない。作用次元を持つシンプレクティック面積と単なる振動数積
\(\nu_1\nu_2\)
を同一視するには本来この写像とスケールの指定が要り、本稿はそれを観察上の便宜として採るにすぎない（この写像も「本稿後半について」で述べた対応仮説に属する）。

\subsubsection{\texorpdfstring{4.2 比 \(k\)
と等方・非等方}{4.2 比 k と等方・非等方}}\label{ux6bd4-k-ux3068ux7b49ux65b9ux975eux7b49ux65b9}

二軸の比を

\[k\equiv\frac{\nu_2}{\nu_1}\]

と定義する。\(k=1\)
は二軸が対等な等方の状態（位相面の上で円）であり、\(k\neq1\)
は非等方（楕円）である。\(k\)
は面積を保つスクイーズのパラメータであり、\(k=1\) が等方、\(k\neq1\)
が非等方の歪みに対応する------以下の議論はこの幾何だけで閉じる。（同種のスクイーズを速度パラメータと結ぶ対応は同著者
{[}10{]} でも扱うが、本節はその対応に依存しない。）

非等方が生じたとき、二軸のうち一方（比が詰まる側）を「空間方向の振動数（波長側）」、他方を「時間方向の振動数」と後から名づけることができる。この見方では、空間と時間の区別を、二軸に先験的に刻まれた性質ではなく、非等方が生むラベルとして読むことができる。ただしこれは本稿の読み方であって、標準理論における時空構造の決定機構を置き換える主張ではない（証明もしていない）。この読みのもとでは、等方（\(k=1\)）の段階に、空間方向と呼べる特権的な軸はまだ立っていない。

\begin{center}\rule{0.5\linewidth}{0.5pt}\end{center}

\subsection{§5
従属量としての四元運動量}\label{ux5f93ux5c5eux91cfux3068ux3057ux3066ux306eux56dbux5143ux904bux52d5ux91cf}

\subsubsection{5.1
運動量は振動数の従属量}\label{ux904bux52d5ux91cfux306fux632fux52d5ux6570ux306eux5f93ux5c5eux91cf}

\(p=\hbar k\)（\(k\)
は波数＝空間方向の振動数）であり、運動量は空間方向の振動数に換算定数を掛けた量として書ける（de
Broglie
関係）。本稿はこの関係を用いて、位相面の自由度を、位置と運動量を別々にではなく\textbf{振動数の本数}で数える取り方を採る。これは標準的な位相空間が位置・運動量を独立座標として扱うことを否定するものではなく、同じ自由度を、振動数を主役にして数え直すという本稿の数え方である。時間が従属ラベルであったのと同じ構造を、運動量についても採る。

\subsubsection{5.2
多軸への拡張と四元運動量}\label{ux591aux8ef8ux3078ux306eux62e1ux5f35ux3068ux56dbux5143ux904bux52d5ux91cf}

二軸 \((\nu_1,\nu_2)\) の位相面を多軸へ拡張し、四本の振動数
\(\nu_1,\dots,\nu_4\) を考える（次元が四である理由は §6
で観察する）。四本それぞれに従属量がぶら下がる。四本のうち三本を空間方向、一本を時間方向と（後から）振り分ければ、前者の従属量が運動量の三成分
\(\boldsymbol{p}=(p_x,p_y,p_z)\) に、後者の従属量がエネルギー \(E\)
にあたる。\(E=\hbar\nu\) と \(\boldsymbol{p}=\hbar\boldsymbol{k}\)
は、別々の関係式ではなく、「振動数に従属量がつく」という同一構造の四成分である。ただしここで得られるのは、四つの振動数様変数を並べた\textbf{形式的な四成分量}までである。これが物理的な四元運動量
\(p_\mu=(E,\boldsymbol{p})\) になるには、背後にミンコフスキー計量
\(\eta_{\mu\nu}\)
と時空並進対称性があってローレンツ変換の表現として振る舞う必要がある。本稿は符号構造を
§8.4
のとおり外しているため、現段階で言えるのは「ミンコフスキー計量と並進対称性を別途導入できるならば、この四成分量は四元運動量と対応づけられる可能性がある」までである。

ここで重要な観察を記録する。\textbf{本稿の構成では、この段階まで四本の振動数は対等に扱われており、どれが空間でどれが時間かを分ける情報を本稿はまだ導入していない。}
一本だけ符号が反転する（ミンコフスキー計量）という構造は、この対称な段階で本稿が導入する性質ではなく、後から何かが対称性を破って初めて入るもの、と本稿は見る。何がその対称性を破るかは本稿では扱わず（§8.4
の未解決課題）、本稿は対称な四本までを扱う。これは、ミンコフスキー計量や符号構造が標準理論で後発的だと主張するものではなく、本稿の思考実験がそこまでしか踏み込まない、という範囲の限定である。

\begin{center}\rule{0.5\linewidth}{0.5pt}\end{center}

\subsection{§6
なぜ振動数は四本か------多元体の離散条件}\label{ux306aux305cux632fux52d5ux6570ux306fux56dbux672cux304bux591aux5143ux4f53ux306eux96e2ux6563ux6761ux4ef6}

\subsubsection{6.1
床は次元を整数に固定しない}\label{ux5e8aux306fux6b21ux5143ux3092ux6574ux6570ux306bux56faux5b9aux3057ux306aux3044}

位相面積の床 \(\tfrac12\)
は連続的な制約である。点に潰れられないことは保証するが、それ自体では次元を整数値に固定しない。もし床だけが効いているなら、次元は三と四、四と五の間で定まらずに揺らぐ余地が残る。次元を四という鋭い整数値に釘付けにするには、床よりも厳しい離散的な条件が要る。それが多元体の存在条件である。

\subsubsection{6.2 許される次元は 1, 2, 4, 8
に限られる}\label{ux8a31ux3055ux308cux308bux6b21ux5143ux306f-1-2-4-8-ux306bux9650ux3089ux308cux308b}

二軸が張る面積（ノルム）が、軸の合成のもとで積として保たれる------これが「振動の合成がノルムを壊さない」条件である。この条件を満たす代数（平方和の積公式が成り立つノルム付き多元体）は、実数
\(\mathbb{R}\)・複素数 \(\mathbb{C}\)・四元数 \(\mathbb{H}\)・八元数
\(\mathbb{O}\) の四つに限られ、その次元は 1, 2, 4, 8 である（Hurwitz
1898/1923
{[}5{]}）。三次元のノルム付き多元体は存在しない。ここは限定が要る。三次元を扱う代数構造（回転
\(\mathfrak{so}(3)\)、ベクトル積、Clifford
代数など）は多数あり、「三次元に代数がない」のではない。存在しないのは\textbf{三本の実基底を持つノルム付き除法多元体}であって、「三本では（この意味で）閉じない」という直感は、その定理に限ってのみ裏づけられる。なお、結合性を落とした実多元体一般でも、許される次元は
1, 2, 4, 8 に限られることが代数的位相幾何学（Bott--Milnor、Kervaire 1958
{[}7{]}）によって示されている。

\subsubsection{6.3
ここからは観察ではなく著者の解釈である------四を指す読み方}\label{ux3053ux3053ux304bux3089ux306fux89b3ux5bdfux3067ux306fux306aux304fux8457ux8005ux306eux89e3ux91c8ux3067ux3042ux308bux56dbux3092ux6307ux3059ux8aadux307fux65b9}

\textbf{本節以降は性格が変わることを、最初に明示する。} §6.2
までは、既知の数学------不確定性の床、およびノルム付き多元体の次元が 1,
2, 4, 8
に限られること------を並置して観察した。ここは誰も反論できない既知の構造の並置である。\textbf{これに対し、本節（§6.3）と次節（§7）は、それらの既知構造を著者がどう読むかという解釈・直感であって、観察論文としての主張ではない。}
以下で「四が選ばれる」と述べるとき、それは標準量子論や数学の定理が四を要求するという意味ではない。著者が一定の条件を時空四軸の合成代数の上に置いて読むと四に行き着く、という著者の見立ての記述である。読者はこれを観察結果ではなく、著者の解釈として受け取られたい。

\paragraph{移す操作それ自体が仮定である}\label{ux79fbux3059ux64cdux4f5cux305dux308cux81eaux4f53ux304cux4eeeux5b9aux3067ux3042ux308b}

許される 1, 2, 4, 8
のどれを時空四軸の代数に当てるかを考える。本節が置く条件は、標準量子論の構造に現れる対称性を、この四軸の合成代数の上へ\textbf{移して読んだもの}である。ここを正確に述べる。これらの対称性（ユニタリ性・可逆性・角運動量の非可換性・演算子合成の結合性）は、本来は\textbf{状態空間（複素ヒルベルト空間）と演算子代数の性質}であって、時空次元に対して何かを要求するものではない。標準量子論は二次元（\(\mathbb{C}\)）のヒルベルト空間で平然と成立する。したがって「これらを時空四軸の代数に課す」という\textbf{移す操作それ自体が、本稿の（証明されない）取り方}であり、その正当化は本稿にはなく、著者の直感に属する。結論の全重量はこの一手に乗っている------これを但し書きの隅ではなく本文で引き受ける。

\paragraph{移したあとで効く条件は、実質二つ}\label{ux79fbux3057ux305fux3042ux3068ux3067ux52b9ux304fux6761ux4ef6ux306fux5b9fux8ceaux4e8cux3064}

移す条件を挙げると、独立に効くのは次の二つである。

\begin{itemize}
\tightlist
\item
  \textbf{積が非可換であること}：可換な代数（\(\mathbb{R}, \mathbb{C}\)）では独立した二枚の回転面が立たない。これにより次元
  1, 2 を落とす。
\item
  \textbf{積が結合的であること}：合成順序が定義できること。八元数
  \(\mathbb{O}\) は非結合なので、これにより次元 8 を落とす。
\end{itemize}

「ノルム保存（ユニタリ性）」と「可逆性（除法可能）」も条件として挙げうるが、ノルム付き多元体は同時に除法多元体でもあるため、両者は同じ次元集合
\(\{1,2,4,8\}\) を選び、独立な絞り込みを足さない（Hurwitz
{[}5{]}）。すなわち \(\{1,2,4,8\}\)
の中で四を指す独立条件は、\textbf{非可換性と結合性の二つだけ}である。表にすると：

{\def\LTcaptype{none} % do not increment counter
\begin{longtable}[]{@{}
  >{\raggedright\arraybackslash}p{(\linewidth - 10\tabcolsep) * \real{0.1154}}
  >{\raggedright\arraybackslash}p{(\linewidth - 10\tabcolsep) * \real{0.1154}}
  >{\centering\arraybackslash}p{(\linewidth - 10\tabcolsep) * \real{0.1923}}
  >{\centering\arraybackslash}p{(\linewidth - 10\tabcolsep) * \real{0.1923}}
  >{\centering\arraybackslash}p{(\linewidth - 10\tabcolsep) * \real{0.1923}}
  >{\centering\arraybackslash}p{(\linewidth - 10\tabcolsep) * \real{0.1923}}@{}}
\toprule\noalign{}
\begin{minipage}[b]{\linewidth}\raggedright
次元
\end{minipage} & \begin{minipage}[b]{\linewidth}\raggedright
代数
\end{minipage} & \begin{minipage}[b]{\linewidth}\centering
ノルム付き多元体（＝可除）
\end{minipage} & \begin{minipage}[b]{\linewidth}\centering
積が非可換
\end{minipage} & \begin{minipage}[b]{\linewidth}\centering
積が結合的
\end{minipage} & \begin{minipage}[b]{\linewidth}\centering
全条件
\end{minipage} \\
\midrule\noalign{}
\endhead
\bottomrule\noalign{}
\endlastfoot
1 & \(\mathbb{R}\) & $\bigcirc$ & $\times$ & $\bigcirc$ & $\times$ \\
2 & \(\mathbb{C}\) & $\bigcirc$ & $\times$ & $\bigcirc$ & $\times$ \\
3 & --- & $\times$ & --- & --- & $\times$ \\
\textbf{4} & \(\mathbb{H}\) & $\bigcirc$ & $\bigcirc$ & $\bigcirc$ & \textbf{$\bigcirc$} \\
5 & --- & $\times$ & --- & --- & $\times$ \\
6 & --- & $\times$ & --- & --- & $\times$ \\
8 & \(\mathbb{O}\) & $\bigcirc$ & $\bigcirc$ & $\times$ & $\times$ \\
\end{longtable}
}

（最初の列はすべて「ノルム付き＝可除」を一つにまとめた。3, 5, 6
はそもそもノルム付き多元体が存在しない。）この読みのもとで、可除・非可換・結合を同時に満たす最小次元は四（四元数
\(\mathbb{H}\)）である（結合的実多元体は同型を除き
\(\mathbb{R},\mathbb{C},\mathbb{H}\) の三つ：Frobenius
{[}6{]}）。離散構造との整合（Hurwitz 整数・\(D_4\) 格子）は §6.4。

\paragraph{循環についての自己批判}\label{ux5faaux74b0ux306bux3064ux3044ux3066ux306eux81eaux5df1ux6279ux5224}

上で「非可換性」を、角運動量代数
\([J_i,J_j]=i\varepsilon_{ijk}J_k\)（\(\mathrm{SU}(2)\)）の非可換性として動機づけると、その非可換性は三次元空間の回転
\(\mathrm{SO}(3)\)
に由来する。三次元空間を前提して四を導くことになり、\textbf{循環}である。これを避けるには、ここで効かせる非可換性を\textbf{四元数の積の非可換性}に純化しなければならない。ただし角運動量代数の非可換性と
\(\mathbb{H}\)
の積の非可換性は構造的に類似するが\textbf{同一ではなく}、両者を同一視すること自体もまた本稿の（証明されない）見立てである。この点でも、本節は導出ではなく著者の読み方にとどまる。

\paragraph{標準量子論との関係}\label{ux6a19ux6e96ux91cfux5b50ux8ad6ux3068ux306eux95a2ux4fc2}

念のため繰り返す。状態の重ね合わせと確率振幅が複素数で書かれること（複素ヒルベルト空間）は標準どおり継承し、四元数と結びつくのは状態のスカラー体ではなく、本稿が条件を移した時空四軸の合成代数である。本節は標準量子論のいかなる予言とも矛盾しない。系が一次元から四へ拡張する過程、五次元以上へ進まない理由は、§7
で著者の推測として扱う。

\subsubsection{6.4
離散整数としての四本}\label{ux96e2ux6563ux6574ux6570ux3068ux3057ux3066ux306eux56dbux672c}

四元数の整数版には二種類あり、ここは区別が要る。Lipschitz
整数（座標がすべて整数の四元数）は \(\mathbb{Z}^4\)
格子と自然に対応する。Hurwitz
整数はそれに半整数座標の点を加えた、より対称性の高い四次元格子構造であり、これは
\(D_4\) 格子と密接に関連する。すなわち \(D_4\)・Hurwitz
整数・\(\mathbb{Z}^4\)
は互いに関連するが\textbf{単純に同一の対象ではない}（\(D_4\) を単なる
\(\mathbb{Z}^4\) と同一視すべきではない）。\(D_4\)
格子の最密充填の一意性は古典的に知られている
{[}9{]}。複素数（ガウス整数）が二次元の離散版であるのに対し、回転が非可換かつ結合的で、整数格子の上で全方向が噛み合う向きとして四が現れる。四本の振動数
\(\nu_1,\dots,\nu_4\) は、これらの離散整数の上に乗る量として読める。表記
\(\pm\tfrac12\) は、整数を点に潰させない最小の連続的余地（整数 \(n\)
＋絞りきれない半端な床、すなわち \(n+\tfrac12\) の
\(\tfrac12\)）として、四本それぞれに効く。

\subsubsection{6.5
物理側の独立な系列との符合（参考）}\label{ux7269ux7406ux5074ux306eux72ecux7acbux306aux7cfbux5217ux3068ux306eux7b26ux5408ux53c2ux8003}

参考として、これは物理側の独立な定理とも噛み合う。超対称な Yang--Mills
理論が成立する時空次元が 3, 4, 6, 10 に限られ、その理由が二つ下の次元 1,
2, 4, 8
のノルム付き多元体の存在にある、という対応が知られている（Baez--Huerta
{[}8{]}）。本稿はこの対応を導出も主張もしないが、四本の振動数が四元数に乗るという見立てが、既知の系列の上に素直に位置することを記録する。

\begin{center}\rule{0.5\linewidth}{0.5pt}\end{center}

\subsection{§7
四で安定する理由についての推測}\label{ux56dbux3067ux5b89ux5b9aux3059ux308bux7406ux7531ux306bux3064ux3044ux3066ux306eux63a8ux6e2c}

\textbf{本節は証明ではなく、推測である。} §6.3
までで「標準量子論の対称性を時空四軸の代数へ移して読むと、それらを満たす最小次元は四である」という（対応仮説のもとでの）静的な見立ては確認した。残るのは、§6.3
末尾と §8.4
で後段送りとした問い------系がなぜ四まで達し、四で止まるのか------である。本節はこれに対し、確定した答えではなく、一つの推測を観測として記すに留める。

\subsubsection{7.1
仮定：次元数もまた不確定性を負う}\label{ux4eeeux5b9aux6b21ux5143ux6570ux3082ux307eux305fux4e0dux78baux5b9aux6027ux3092ux8ca0ux3046}

本稿はここで一つの仮定を明示的に置く（導出ではない）。位置・運動量が共役対として不確定性の床を負い、面積ゼロの点状状態に潰れられない------これとの\textbf{類比として}、\textbf{次元数もまた揺らぐ量であり、特別な対称性の破れがない限り、確定した整数ではなく揺らぎを持つ}ものと仮定する。

ここで地位の違いを明確に切り分ける。§2 の床は、正準共役対 \([q,p]=i\)
から Robertson
の不等式として厳密に出る。これに対し本節の「次元数の不確定性」は、次元数
\(D\) に共役な観測量が何かを本稿が指定していない以上、Robertson
型の厳密な床ではなく\textbf{類比}である。これを厳密にするには、次元数演算子・その共役量・次元揺らぎを載せる状態空間・次元間遷移の力学といった追加構造が要るが、本稿はそれを与えない。したがって
§2
の床と本節の「次元の揺らぎ」は地位が異なり、後者はあくまで著者の推測の出発点である。

次元を固定された整数ではなく動的・スケール依存の量として扱う見方には、類似する先行研究が存在する。量子重力のいくつかの独立なアプローチが短距離で有効次元の変化（次元的縮約）を示すこと
{[}12{]}、因果力学的単体分割でスペクトル次元がスケール依存に走ること
{[}13{]} が知られ、とくに {[}12{]}
はこの縮約を場の理論の自発的対称性の破れに類比して論じている。\textbf{ただし、これらは次元を非確定の量として扱う見方が突飛でないことを示す文脈であって、本稿の仮定の論拠ではない。}
方向もむしろ逆（これらは短距離で次元が減る向きを扱う）であり、本稿はこれらの結果に依存しない。上の仮定はあくまで本稿が前提として置くものであり、その帰結を以下に観測する。

\subsubsection{7.2
推測：揺らぎによる拡大と四での安定化}\label{ux63a8ux6e2cux63faux3089ux304eux306bux3088ux308bux62e1ux5927ux3068ux56dbux3067ux306eux5b89ux5b9aux5316}

この仮定のもとで、次の推測を順に置く。いずれも本稿は定量化の手段を持たず、確率を明示できない。

第一に、通常の意味での零次元には正準位相面が存在しないため、§2 の
Robertson の床をそのまま適用することはできない。ここは §7.1
の類比として置く------次元数にも揺らぎを許すならば、「零次元」は完全に確定した無ではなく、次元生成の可能性を含む未確定の状態として読める、と仮定する。

第二に、この揺らぎにより、零次元から一次元以上へ次元が拡大する確率が存在すると仮定する。確率の値も拡大の機構も、本稿は与えない。

第三に、次元は安定しないまま揺らぎ、四次元まで拡大しうると仮定する。

第四に、いったん四次元に達すると、§6.3
で見た多元体条件------標準量子論の対称性を移して読んだときに四軸の代数が満たすべき条件------が働く。これは連続的な不確定性の床よりも強い、離散的な拘束である。四次元はこの拘束を満たす最小の次元であり、ここで系は安定状態に落ち着く。

第五に、この安定性は不確定性の揺らぎよりも強いため、揺らぎはその障壁を越えられず、四次元から低次元・高次元へ再び揺らぐことが禁止される------のではないか。すなわち四は、「揺らぎで到達でき、かつ到達すると揺らぎでは出られない」次元として選ばれる。

繰り返すが、第四・第五の「四で安定し、障壁を越えられない」という描像そのものも、本稿が証明するものではなく、推測として置くものである。とくに第五は推測群の中でも飛躍が最大であることを、自覚的に記す。§6.3
の多元体条件は\textbf{静的な代数的事実}（四でのみ可除・非可換・結合が揃う）であって、「閉じる」ことと、力学的に「安定点で障壁を越えられない」こととは別物である。後者を言うにはポテンシャル・作用といった動力学が要るが、本稿はそれを持たない。ポテンシャルの谷・安定の障壁という比喩も、定量化はしない。評価は読者に委ねる。

\subsubsection{7.3
二制約の合流（記録に留める）}\label{ux4e8cux5236ux7d04ux306eux5408ux6d41ux8a18ux9332ux306bux7559ux3081ux308b}

この推測を整理すると、次元四は二つの独立な制約の合流点として現れる。一つは連続的な床（点になれない、ゆえに次元が零に潰れず揺らぎうる）。もう一つは離散的な代数条件（揺らぎが四に達すると、そこで初めて閉じ、四を安定点とする）。前者は不確定性、後者は多元体という、出自の異なる二つの数学が、ともに同じ次元四を指す。

本稿はこの符合を\textbf{記録する}に留め、これを偶然と見るか必然と見るか、あるいはより深い単一の根に還元できるかについては、評価を読者に委ねる。本稿の枠組みからは、この二制約のいずれもが導出されるのではなく、外部の確立された数学として引用されている。

\begin{center}\rule{0.5\linewidth}{0.5pt}\end{center}

\subsection{§8
既存研究との関係と、本稿が主張しないこと}\label{ux65e2ux5b58ux7814ux7a76ux3068ux306eux95a2ux4fc2ux3068ux672cux7a3fux304cux4e3bux5f35ux3057ux306aux3044ux3053ux3068}

\subsubsection{8.1
既存研究との関係}\label{ux65e2ux5b58ux7814ux7a76ux3068ux306eux95a2ux4fc2}

本稿の各構成要素は標準的に確立されている。位相空間の面積による不確定性の解釈は
Robertson 1929 {[}2{]}、時間--帯域版は Gabor 1946
{[}3{]}、面積の不変量としての量子ブロブは de Gosson
{[}4{]}。ノルム付き多元体の次元が 1, 2, 4, 8 に限られることは Hurwitz
{[}5{]}、結合的実多元体が \(\mathbb{R},\mathbb{C},\mathbb{H}\)
に限られることは Frobenius {[}6{]}、実多元体一般の次元制限は
Bott--Milnor / Kervaire {[}7{]}、多元体と超対称性の対応は Baez--Huerta
{[}8{]}、\(D_4\) 格子の充填は
{[}9{]}。位相空間の面積を共通通貨とする視点、およびスクイーズ＝ブースト対応は同著者観察論文
{[}10{]}。本稿はこれらを、一つのモデル粒子の思考実験という設定の中で、位相面積の視点から並べ直したに過ぎない。

\subsubsection{8.2
本稿が主張しないこと}\label{ux672cux7a3fux304cux4e3bux5f35ux3057ux306aux3044ux3053ux3068}

\begin{itemize}
\tightlist
\item
  標準量子論・場の量子論・特殊相対論の数学的予言の変更（本稿はこれらと観測上等価である）。
\item
  時空次元が 3+1
  であることの物理的導出。本稿が観察するのは、ノルムが積で閉じる代数の許容次元という数学的条件であって、現実の時空次元の決定機構ではない。両者を同一視しない。
\item
  計量符号・ミンコフスキー構造・虚数単位の導出。四本の振動数が対等な段階（§5.2）までを扱い、符号の反転は導入しない。
\item
  新しい物理定数・断面積・崩壊率・新粒子の予言。
\item
  新しい数学的定理の証明。引用した定理（Hurwitz、Frobenius
  等）はすべて外部の確立された結果である。
\item
  宇宙の起源・特異点・「設計」に関する形而上学的主張。§7
  は「次元数もまた不確定性を負う」という仮定のもとでの、四での安定化についての推測を述べるに留まり（証明ではない）、それ以上の解釈は読者に委ねる。
\end{itemize}

\subsubsection{8.3
本稿が記録すること}\label{ux672cux7a3fux304cux8a18ux9332ux3059ux308bux3053ux3068}

\begin{itemize}
\tightlist
\item
  モデル粒子は位相面積の床 \(\tfrac12\)
  ゆえに点状状態を取れず、必ず有限の広がりを持つこと。等方最小不確定状態の各軸標準偏差は
  \(\tfrac{1}{\sqrt2}\) であり、本稿の表記 \(\pm\tfrac12\)
  は各軸の幅ではなく床 \(\tfrac12\)（ゼロ点の半端）そのものを指すこと。
\item
  位相面の二軸を振幅と振動数に取り、ともに振動数の次元へ無次元化すると、面積が純粋な積
  \(\nu_1\cdot\nu_2\)
  となること。この枠では時間を従属ラベルとして扱い、面の主軸を振幅と振動数に取る取り方が矛盾なく採れること（時間を主軸に据える取り方を排除はしない）。
\item
  比 \(k=\nu_2/\nu_1\)
  が等方（\(k=1\)）と非等方（\(k\neq1\)）を分け、空間・時間の区別を、二軸の非等方が生むラベルとして読めること（標準理論の時空構造を置き換える主張ではない）。
\item
  運動量とエネルギーを空間方向・時間方向の振動数の従属量として読むと、四本の振動数の従属量を並べた形式的な四成分量が得られること。これが物理的な四元運動量になるにはミンコフスキー計量の別途導入が要り、四本の段階では符号の非対称が未導入であること。
\item
  次元を四に固定するのが、連続の床ではなく、多元体の離散条件であること。標準量子論の構造に現れる対称性（ノルム保存・除法可能・回転の非可換・合成の結合）を振動数四軸の合成代数へ移して読むと、それらを同時に満たす最小次元が四であること。これらの対称性は本稿が後付けで選んだものではなく標準量子論に実在する構造だが、それを時空四軸の代数へ移すこと自体は本稿の取り方であり、状態空間が複素数体であることは標準どおり継承すること（§6.3
  の一覧表と但し書き）。
\item
  「次元数もまた不確定性を負う」と仮定すると、四は「揺らぎで到達でき、到達すると揺らぎでは出られない」安定点として選ばれうること（§7、推測。証明ではなく、類似する先行研究
  {[}12{]}{[}13{]} は論拠ではない）。
\item
  連続の床と離散の環という二つの独立な制約が、ともに次元四に合流するという符合（解釈は読者に委ねる）。
\end{itemize}

\subsubsection{8.4
今後の課題（未解決として明示する点）}\label{ux4ecaux5f8cux306eux8ab2ux984cux672aux89e3ux6c7aux3068ux3057ux3066ux660eux793aux3059ux308bux70b9}

\begin{enumerate}
\def\labelenumi{\arabic{enumi}.}
\tightlist
\item
  \textbf{次元拡張の動的機構}：§6.3
  が明示したのは「標準量子論の対称性を時空四軸の代数へ移して読むと、それらを満たす最小次元は四である」という（対応仮説のもとでの）静的な見立てまでである。系が一次元から四へと対称性を積み増しながら拡張する過程、および五次元以上へ進まない理由は、この事実とは別の問いであり、本稿では扱わない。拡張が動的な発展なのか、それとも成立条件の論理的順序にすぎないのかを含めて、別稿に委ねる。
\item
  \textbf{符号（ミンコフスキー計量）の導入機構}：四本の対等な振動数のうち一本が双曲側（時間）に落ちる機構は本稿では扱わない。何が対称性を破るのか（観測者の射影の選択か、内在の非対称か）は未解決であり、別稿に委ねる。
\item
  \textbf{現実の時空次元との接続}：本稿の「許容次元としての四」と、観測される時空次元
  3+1 の関係（空間三・時間一の内訳がどう決まるか）は、§5.2
  の符号問題と同じく未解決である。
\item
  \textbf{低次元トイモデルの明示化}：二軸 \((\nu_1,\nu_2)\)
  から四本への拡張を、四元数 \(\mathbb{H}\)
  の具体的な作用として書き下し、床 \(\tfrac12\)
  が各軸にどう分配されるかを陽に示すこと。
\item
  \textbf{八次元位相空間との橋渡し}：標準的なシンプレクティック構造では、時空各軸
  \(x^\mu\) に運動量 \(p_\mu\) が共役で、位相空間は八次元である。§5.1
  のように運動量を振動数の従属量として畳むと、この正準対構造が見えなくなる。「振動数の本数で数える」本稿の数え方と八次元位相空間との対応づけは本稿では与えておらず、未解決として残す。
\item
  \textbf{同著者枠組みとの統合}：本稿の振動数四本と、同著者の
  \(\mathbb{Z}^4\) 離散時空・\(D_4\)
  格子・標準模型対応の枠組みとの整合を、別稿で詳細に詰めること。
\end{enumerate}

\begin{center}\rule{0.5\linewidth}{0.5pt}\end{center}

\subsection{関連研究（同著者）}\label{ux95a2ux9023ux7814ux7a76ux540cux8457ux8005}

本稿は外部の標準文献のみで自己完結する。同著者による観察論文
{[}10{]}（位相空間の面積を共通通貨とするシンプレクティック対称性・Stone
の定理・Wick 回転の再読）および
{[}11{]}（信号論・制御理論と量子論の構造的対応の観察）が同主題の周辺にあるが、本稿はそれらの結論に依存しない。本稿の思考実験は、{[}10{]}
の位相面積の枠組みを前提として参照するに留まる。

\begin{center}\rule{0.5\linewidth}{0.5pt}\end{center}

\subsection{参考文献}\label{ux53c2ux8003ux6587ux732e}

{[}1{]} W. Heisenberg (1927). \emph{Über den anschaulichen Inhalt der
quantentheoretischen Kinematik und Mechanik}. Z. Phys. \textbf{43}, 172.
{[}2{]} H. P. Robertson (1929). \emph{The uncertainty principle}. Phys.
Rev.~\textbf{34}, 163. {[}3{]} D. Gabor (1946). \emph{Theory of
communication}. J. IEE \textbf{93}, 429. {[}4{]} M. de Gosson (2013).
\emph{Quantum blobs}. Found. Phys. \textbf{43}, 440. {[}5{]} A. Hurwitz
(1898). \emph{Über die Composition der quadratischen Formen von beliebig
vielen Variablen}. Nachr. Ges. Wiss. Göttingen, 309；(1923) Math. Ann.
\textbf{88}, 1.（ノルム付き多元体の次元が 1, 2, 4, 8 に限られること）
{[}6{]} F. G. Frobenius (1878). \emph{Über lineare Substitutionen und
bilineare Formen}. J. Reine Angew. Math. \textbf{84},
1.（結合的実多元体は \(\mathbb{R},\mathbb{C},\mathbb{H}\) に限られる）
{[}7{]} R. Bott, J. Milnor (1958). \emph{On the parallelizability of the
spheres}. Bull. Amer. Math. Soc. \textbf{64}, 87；M. Kervaire (1958).
\emph{Non-parallelizability of the n-sphere for n\textgreater7}. Proc.
Natl. Acad. Sci. \textbf{44}, 280.（実多元体の次元は 1, 2, 4, 8
に限られる） {[}8{]} J. C. Baez, J. Huerta (2010). \emph{Division
algebras and supersymmetry I}. In \emph{Superstrings, Geometry,
Topology, and C*-algebras}, Proc. Sympos. Pure Math. \textbf{81},
65；arXiv:0909.0551. {[}9{]} J. H. Conway, N. J. A. Sloane (1999).
\emph{Sphere Packings, Lattices and Groups}, 3rd ed.~Springer.（\(D_4\)
格子） {[}10{]} N. Kihara (2026).
\emph{位相空間の面積としての保存量と不確定性------シンプレクティック対称性・Wick
回転・Stone の定理の統一的視点}. 観察論文（同著者）. Concept DOI:
10.5281/zenodo.20521566. {[}11{]} N. Kihara (2026).
\emph{信号論・制御理論と量子力学・量子光学・開放量子系の構造的対応関係に関する観察}.
観察論文（同著者）. Concept DOI: 10.5281/zenodo.20521598. {[}12{]} S.
Carlip (2017). \emph{Dimension and dimensional reduction in quantum
gravity}. Class. Quantum Grav. \textbf{34},
193001；arXiv:1705.05417.（次元を非確定の有効量として扱う見方・自発的対称性の破れとの類比。本稿
§7 の仮定の論拠ではなく、類似する文脈として参照） {[}13{]} J. Ambjørn,
J. Jurkiewicz, R. Loll (2005). \emph{The spectral dimension of the
universe is scale dependent}. Phys. Rev.~Lett. \textbf{95},
171301；arXiv:hep-th/0505113.（スペクトル次元のスケール依存。同上、論拠ではなく文脈として参照）

\begin{center}\rule{0.5\linewidth}{0.5pt}\end{center}

著者：木原 範昭（Noriaki Kihara）／ WF System Co., Ltd.~／ ORCID
\href{https://orcid.org/0009-0004-6753-4020}{0009-0004-6753-4020}／ CC
BY 4.0

\begin{center}\rule{0.5\linewidth}{0.5pt}\end{center}

\subsection{改訂履歴}\label{ux6539ux8a02ux5c65ux6b74}

\begin{itemize}
\tightlist
\item
  \textbf{v1（2026-06）}：初版（ステルス・ドラフト）。一つのモデル粒子の思考実験として、(§1)
  設定、(§2) 床 \(\tfrac12\) と位相面の内在・点状状態の不在、(§3)
  振幅と振動数の二軸・従属ラベルとしての時間、(§4) 無次元化と比
  \(k\)・等方/非等方、(§5) 従属量としての四成分量・四本の対等性、(§6)
  多元体の離散条件による次元四の許容（Hurwitz/Frobenius、標準量子論の対称性を時空四軸の代数へ移して読む一覧表で、全充足の最小次元が四であることを明示。拡張の動的機構は後段送り）、(§7)
  「次元数もまた不確定性を負う」という仮定のもとで、零次元から四への揺らぎによる拡大と四での安定化を推測として観測（証明ではない／類似する先行研究
  Carlip {[}12{]}・AJL {[}13{]}
  は論拠ではなく文脈）、および二制約の合流の記録、を並置観察。符号（ミンコフスキー計量）の導入は今後の課題として明示的に外した。同著者観察論文
  {[}10{]}{[}11{]} の枠組みを前提参照。
\item
  \textbf{v1 改訂（査読反映、2026-06）}：3 AI 査読（Claude.ai / ChatGPT
  / Grok）を反映。(a) §2 と §3
  の境界に「本稿後半について」を新設し、既知の数学（第一層）・著者の対応仮説（第二層、§3〜§6.3
  の時空四軸への読み替え）・推測（第三層、§7）の三層を明示、後半は「導出ではなく作業仮説、観察論文としての主張ではない」と宣言。(b)
  §2.3 の \(\pm\tfrac12\)
  を各軸標準偏差（\(1/\sqrt2\)）と切り分け、ゼロ点の半端（床
  \(\tfrac12\) そのもの）と再定義。(c) §6.3
  の表を「ノルム保存＝可除」で一列に畳み、独立に効くのは非可換性と結合性の二つと明示。「QM
  が要求」を「移して読む」に降格し、角運動量非可換性の循環を自己批判。(d)
  §3.1 で振幅・振動数が正準共役対でないこと、§4.1 で \(\hbar\)
  の役割は消えないこと、§5.2
  で四元運動量はミンコフスキー計量なしでは形式的四成分量に留まることを明示。(e)
  §6.2「三本で閉じない」を「三本の実基底のノルム付き除法多元体に限る」と限定、§6.4
  で Lipschitz/Hurwitz/\(\mathbb{Z}^4\)/\(D_4\) の同一視を解消。(f) §7.1
  で次元数の不確定性は共役量未指定ゆえ類比であることを切り分け、§7.2
  で代数的「閉じる」と力学的「安定」の別を自覚。タイトルを「次元の決定へ」→「次元の許容条件へ」に。
\end{itemize}

\end{document}
